%%%%%%%%%%%%%%%%%%%%%%%%%%%%%%%%%%%%%%%%%%%%%%%%%%%%%%%%%%%%%%%%%%%%%%%%%%%
%
% Generic template for TFC/TFM/TFG/Tesis
%
% By:
%  + Javier Macías-Guarasa.
%    Departamento de Electrónica
%    Universidad de Alcalá
%  + Roberto Barra-Chicote.
%    Departamento de Ingeniería Electrónica
%    Universidad Politécnica de Madrid
%
% By: + Javier Macías-Guarasa. Departamento de Electrónica Universidad de Alcalá + Roberto Barra-Chicote. Departamento de Ingeniería Electrónica Universidad Politécnica de Madrid
% 
% Based on original sources by Roberto Barra, Manuel Ocaña, Jesús Nuevo, Pedro Revenga, Fernando Herránz and Noelia Hernández. Thanks a lot to all of them, and to the many anonymous contributors found (thanks to google) that provided help in setting all this up.
%
% See also the additionalContributors.txt file to check the name of additional contributors to this work.
%
% If you think you can add pieces of relevant/useful examples, improvements, please contact us at (macias@depeca.uah.es)
%
% You can freely use this template and please contribute with comments or suggestions!!!
%
%%%%%%%%%%%%%%%%%%%%%%%%%%%%%%%%%%%%%%%%%%%%%%%%%%%%%%%%%%%%%%%%%%%%%%%%%%%

\chapter{Conclusiones y líneas futuras}
\label{cha:concl-y-line}

\section{Conclusiones}

\subsection{Cumplimiento de objetivos}

El objetivo principal de este trabajo ha sido comparar el rendimiento de distintas familias de modelos de procesamiento del lenguaje natural para la detección automática de trastornos de salud mental en español, utilizando el corpus MentalRiskES como marco experimental común. En concreto, se evaluaron modelos BERT especializados mediante \textit{fine-tuning}, LLMs open-source y LLMs propietarios de última generación sobre cuatro tareas: ansiedad, depresión, TCA y clasificación multiclase.

Además del rendimiento predictivo, se incorporaron dimensiones complementarias como la estabilidad inferencial mediante TARa@5, la comparación estadística mediante bootstrap pareado y el análisis rendimiento-coste mediante frontera de Pareto. Esto ha permitido realizar una evaluación más completa que la simple comparación de métricas agregadas.

Los resultados obtenidos permiten afirmar que los objetivos planteados se han cumplido satisfactoriamente. El estudio no solo identifica qué modelos alcanzan mejores resultados en cada tarea, sino también en qué escenarios dichas diferencias son realmente relevantes desde un punto de vista estadístico y práctico.

\subsection{Principales hallazgos}

Los resultados muestran que los LLMs suelen alcanzar el mejor rendimiento absoluto, aunque no existe un modelo universalmente superior. GPT-OSS-20B obtiene el mejor resultado en ansiedad, Gemini 3.1 Pro lidera en depresión y clasificación multiclase, mientras que GPT-5.2 destaca en TCA. Sin embargo, estas ventajas no siempre son amplias ni estadísticamente significativas.

Dentro de los modelos BERT, se observa además una tendencia clara: los modelos preentrenados específicamente para español, como RoBERTa-large-BNE y DistilBETO, presentan un rendimiento superior al de alternativas multilingües como mDeBERTa o Longformer. Esto sugiere que la especialización lingüística sigue siendo un factor muy relevante en tareas clínicas sensibles al matiz semántico y al contexto cultural del lenguaje.

Uno de los hallazgos más relevantes del trabajo es el excelente comportamiento de RoBERTa-large-BNE, que se consolida como el mejor modelo encoder tradicional y compite de forma muy cercana con varios LLMs considerablemente más costosos. En tareas como TCA y clasificación multiclase incluso supera a la mayoría de modelos generativos, confirmando que un buen modelo BERT ajustado específicamente puede seguir siendo una alternativa altamente competitiva.

En contraste, resulta especialmente llamativo el caso de Longformer. Su inclusión en el estudio partía de la hipótesis de que la capacidad de procesar secuencias mucho más largas permitiría capturar mejor patrones contextuales complejos y mejorar el rendimiento predictivo. Sin embargo, los resultados muestran el efecto contrario: Longformer se posiciona sistemáticamente como el peor modelo BERT e incluso presenta comportamientos degenerados en varias tareas binarias, llegando a predecir únicamente la clase mayoritaria. Una explicación probable es que se trata de una arquitectura de gran tamaño que requiere una mayor cantidad de datos para entrenarse de forma eficaz, mientras que el tamaño relativamente reducido del dataset utilizado limita su capacidad de aprendizaje.

Asimismo, GPT-OSS-20B destaca como el LLM open-source más sólido del estudio, especialmente por su equilibrio entre rendimiento y coste. Frente a ello, algunos modelos propietarios como Claude Sonnet 4.6 muestran una eficiencia considerablemente menor en determinados escenarios, particularmente en clasificación multiclase, donde su elevado coste no se traduce en una mejora proporcional del rendimiento.

Por otro lado, las comparaciones estadísticas muestran que muchas diferencias observadas en F1 macro no alcanzan significancia tras la corrección por comparaciones múltiples, especialmente en depresión, TCA y multiclase. Esto no implica necesariamente que dichas diferencias no existan, sino que, con el reducido tamaño del conjunto de evaluación disponible, no puede asegurarse con suficiente evidencia que esa superioridad sea consistente y no esté influida por la variabilidad muestral. Por ello, pequeñas variaciones en los rankings deben interpretarse con cautela y refuerzan la necesidad de complementar las métricas tradicionales con análisis inferenciales más robustos.

\subsection{Implicaciones prácticas}

Desde una perspectiva aplicada, los resultados sugieren que la elección del modelo no debe basarse únicamente en el máximo rendimiento absoluto, sino en el equilibrio entre precisión, coste computacional y estabilidad predictiva.

Si el objetivo es maximizar el rendimiento sin restricciones presupuestarias, modelos como Gemini 3.1 Pro o GPT-5.2 representan las mejores opciones según la tarea. Sin embargo, cuando se considera la viabilidad de despliegue real, RoBERTa-large-BNE emerge como una de las alternativas más recomendables por su excelente relación coste-rendimiento y su robustez estadística. De forma similar, GPT-OSS-20B constituye la opción open-source más competitiva cuando se busca combinar rendimiento elevado con menor dependencia de APIs propietarias.

En conjunto, este trabajo muestra que los LLMs no sustituyen automáticamente a los modelos encoder tradicionales. En muchos escenarios, especialmente cuando se consideran restricciones económicas y reproducibilidad experimental, modelos encoder ajustados específicamente para la tarea mediante fine-tuning, como RoBERTa-large-BNE, continúan ofreciendo una solución más eficiente y práctica para sistemas reales de detección temprana en salud mental.

\section{Líneas futuras}
\label{sec:lineas-futuras}

La principal limitación de este trabajo ha sido el tamaño y la naturaleza del dataset utilizado. Aunque MentalRiskES constituye uno de los recursos más relevantes disponibles para esta tarea en español, el número de ejemplos sigue siendo relativamente reducido, lo que limita tanto la capacidad de generalización de los modelos como la solidez estadística de las comparaciones realizadas. En futuros trabajos sería especialmente recomendable utilizar corpus de mayor tamaño y, preferiblemente, anotados o validados por profesionales de la salud mental, lo que permitiría reforzar el rigor clínico del estudio. Un mayor volumen de datos facilitaría además un mejor ajuste de los modelos encoder mediante fine-tuning, pudiendo incrementar todavía más la competitividad de modelos como RoBERTa-large-BNE frente a los LLMs, además de aportar mayor significancia estadística a las comparaciones entre modelos.

Dentro del uso de LLMs, una línea especialmente interesante consiste en profundizar en estrategias avanzadas de prompting. En este trabajo se utilizó una formulación fija de las instrucciones de clasificación, pero futuras investigaciones podrían analizar cómo distintas variantes de prompt afectan al rendimiento final. Técnicas como few-shot prompting, self-consistency o estrategias más elaboradas de prompt engineering podrían mejorar la capacidad de razonamiento clínico del modelo sin necesidad de reentrenamiento adicional.

Otra línea de gran interés sería la incorporación de sistemas RAG (Retrieval-Augmented Generation) apoyados en bases de conocimiento clínicas especializadas. La posibilidad de recuperar documentos relevantes a partir del texto del usuario (por ejemplo, criterios diagnósticos oficiales, guías clínicas o literatura médica contrastada) permitiría proporcionar contexto adicional al LLM antes de realizar la predicción, favoreciendo decisiones más precisas, explicables y alineadas con criterios clínicos reales.

Finalmente, dado el rápido avance del área, resulta necesario continuar evaluando nuevos LLMs que vayan apareciendo, tanto open-source como propietarios, ya que pequeñas mejoras arquitectónicas o nuevos procesos de alineamiento pueden modificar de forma importante el equilibrio entre rendimiento, coste y estabilidad observado en este trabajo.



%%% Local Variables:
%%% TeX-master: "../book"
%%% End:


